\subsection{Integration of a vector-valued function with respect to a vectorial measure}

PROPOSITION 11. --- Let F, G, H be three Banach spaces, \(\Phi\) a continuous bilinear mapping of F × G into H. Let m be a majorizable vectorial measure on T, with values in G. Then there exists one and only one continuous linear mapping I \(_{\Phi,\mathbf{m}}\) of \(\overline{\mathcal{L}}_{\mathrm{F}}^{1}(|\mathbf{m}|)\) into H such that, for every z ∈ F and every numerical function h integrable for \textbar m\textbar, one has I \(_{\Phi,\mathbf{m}}(hz) = \Phi(z, \int h dm)\) . Moreover,

\[
\left| \mathrm{I} _ {\Phi , \mathbf {m}} (\mathbf {f}) \right| \leqslant \| \Phi \| \int | \mathbf {f} | d | \mathbf {m} |\tag{4}
\]

for every function \(\mathbf{f} \in \overline{\mathcal{L}}_{\mathrm{F}}^{1}(|\mathbf{m}|)\) .

If there exists such a mapping, its value for a step function \(\mathbf{f}\) over the \(|\mathbf{m}|\) -integrable sets is uniquely determined: for, it is known that one can then write \(\mathbf{f} = \sum_{i} \mathbf{a}_{i} \varphi_{\mathrm{X}_{i}}\) , where the \(\mathrm{X}_{i}\) are \(|\mathbf{m}|\) -integrable and disjoint, and the \(\mathbf{a}_{i} \in \mathrm{F}\) (Ch. IV, §4, No.~9, Lemma). The value of \(\mathrm{I}_{\Phi, \mathbf{m}}(\mathbf{f})\) must therefore be equal to \(\sum_{i} \Phi(\mathbf{a}_{i}, \int \varphi_{\mathrm{X}_{i}} d\mathbf{m})\) . Now, we have (No.~3, Prop. 5)

\[
\left| \sum_ {i} \Phi (\mathbf {a} _ {i}, \int \varphi_ {\mathrm{X} _ {i}} d \mathbf {m}) \right| \leqslant \| \Phi \| \cdot \sum_ {i} | \mathbf {a} _ {i} | \cdot | \mathbf {m} | (\mathrm{X} _ {i}) = \| \Phi \| \int | \mathbf {f} | d | \mathbf {m} |,\tag{5}
\]

which shows first of all that the element \(\sum_{i} \Phi(a_i, \int \varphi_{X_i} d\mathbf{m})\) of H does not depend on the particular expression of \(\mathbf{f}\) in the form \(\sum_{i} a_i \varphi_{X_i}\) , hence that we may denote it by \(I_{\Phi, \mathbf{m}}(\mathbf{f})\) . One verifies immediately that the mapping \(I_{\Phi, \mathbf{m}}\) so defined is linear on the space \(\mathcal{E}_F\) of step functions over the \(|\mathbf{m}|\) -integrable sets: for, it suffices to write two functions \(\mathbf{f}, \mathbf{g}\) of \(\mathcal{E}_F\) in the form \(\mathbf{f} = \sum_{i} a_i \varphi_{X_i}\) and \(\mathbf{g} = \sum_{i} b_i \varphi_{X_i}\) with the same finite family of pairwise disjoint \(|\mathbf{m}|\) -integrable sets \(X_i\) (which is possible by virtue of the Lemma of Ch. IV, §4, No.~9). The inequality (5) then shows that \(I_{\Phi, \mathbf{m}}\) is continuous on \(\mathcal{E}_F\) , and since this subspace is dense in \(\overline{\mathscr{L}}_F^1\) (Ch. IV, §4, No.~10, Cor. 1 of Prop. 19 and Ch. V, §1, No.~3), one deduces from this the existence and uniqueness of \(I_{\Phi,m}\) , as well as the inequality (4).

One says that \(I_{\Phi,\mathbf{m}}(\mathbf{f})\) is the integral of f with respect to m (relative to the bilinear mapping \(\Phi\) ); when the value of the bilinear mapping \(\Phi\) at the point \((\mathbf{x},\mathbf{y})\) is denoted xy, we shall write \(\int f dm\) instead of \(I_{\Phi,\mathbf{m}}(\mathbf{f})\) .

With the notations of No.~6, the integral \(\int \langle \mathbf{f},\mathbf{g}\rangle d\mu\) is none other than \(\mathrm{I}_{\Phi ,\mathbf{m}}(\mathbf{f})\) with \(\Phi (\mathbf{x},\mathbf{x}^{\prime}) = \langle \mathbf{x},\mathbf{x}^{\prime}\rangle\) and \(\mathbf{m} = \mathbf{g}\cdot \mu\)

COROLLARY. --- If m and \(m'\) are two majorizable measures on T, with values in G, then \(I_{\Phi,m+m'} = I_{\Phi,m} + I_{\Phi,m'}\) and \(I_{\Phi,\lambda m} = \lambda I_{\Phi,m}\) for every scalar \(\lambda\) .

The second assertion is immediate. The first signifies that for every function f that is both \(|m|\) -integrable and \(|m'|\) -integrable,

\[
\mathrm{I} _ {\Phi , \mathbf {m} + \mathbf {m} ^ {\prime}} (\mathbf {f}) = \mathrm{I} _ {\Phi , \mathbf {m}} (\mathbf {f}) + \mathrm{I} _ {\Phi , \mathbf {m} ^ {\prime}} (\mathbf {f}).\tag{6}
\]

The function \(\mathbf{f}\) is \((|\mathbf{m}| + |\mathbf{m}'|)\) -integrable (Ch. V, §2, No.~2, Cor. 1 of Prop. 3), hence a fortiori \((|\mathbf{m} + \mathbf{m}'|)\) -integrable, and the first member of (6) is indeed meaningful. To show (6), it suffices to do so for \(\mathbf{f}\) a step function over the \((|\mathbf{m}| + |\mathbf{m}'|)\) -integrable sets, since the set of these functions is dense in \(\overline{\mathcal{L}}_{\mathrm{F}}^{1}(|\mathbf{m}| + |\mathbf{m}'|)\) and the two members of (6) are continuous in the latter space, by virtue of (4). But for \(\mathbf{f} = \mathbf{a}\varphi_{\mathrm{X}}\) , where \(\mathrm{X}\) is \((|\mathbf{m}| + |\mathbf{m}'|)\) -integrable, the two members of (6) reduce to \(\Phi(\mathbf{a}, \int \varphi_{\mathrm{X}} d\mathbf{m}) + \Phi(\mathbf{a}, \int \varphi_{\mathrm{X}} d\mathbf{m}')\) , whence the corollary.

Remark. --- When \(\mathbf{m}\) is of the form \(\mathbf{b}\mu\) , where \(\mathbf{b} \in G\) and \(\mu\) is a real measure on \(T\) ,

\[
\mathrm{I} _ {\Phi , \mathbf {m}} (\mathbf {f}) = \int \Phi (\mathbf {f} (t), \mathbf {b}) d \mu (t)
\]

for every function \(\mathbf{f} \in \mathcal{L}_{\mathrm{F}}^{1}(\mu)\) , because both members are continuous on this space and coincide when \(\mathbf{f}\) is a step function over the \(|\mu|\) -integrable sets.

